
\subsection*{Cross-validation and performance metrics}

We evaluated model generality and performance using two complementary cross-validation (CV) approaches. Five CV folds were used in all model runs, with stratified sampling using biomes as strata. The first approach, using 'standard' CV, assessed how well models could make predictions in environmental and taxonomic contexts similar to the data that they were trained on. In other words, this evaluated model generalizability from sample to sampled population, also referred to as interpolation accuracy. The folds were generated by sampling at the site level among all studies within a given stratum. This implies that sites from a given study were generally split among several folds, creating overlap in studies, but not sites, between folds. In the BTMs, fold assignment was still done at site level, keeping multiple taxonomic units for a given site in the same fold, if present.

The second approach, based on cross-study validation \cite{bernau2014}, was used to assess how well the models could make predictions in new contexts; specifically, whether model learning was transferable between different source studies. In this case, folds were constructed by sampling at the study level, such that all sites from a given study ended up in a single fold. The only exception was if a study spanned several strata (biomes), in which case it could appear in more than one fold. The cross-study validation approach led to a clear spatial and environmental separation of training and test data, while also reflecting taxonomic and methodological differences between source studies. Although the numbers of studies per fold were roughly equal, one consequence of the cross-study procedure was that folds in some iterations became unbalanced in terms of the number of sites, due to the large spread in the size of different studies. In this study we defined model accuracy as the Pearson correlation coefficient between predicted values $\widehat{\mathbf{y}}$ and observed values $\mathbf{y}$:
\vspace{0.25\baselineskip}
\begin{equation*}
    r(\widehat{\mathbf{y}}, \mathbf{y}) = \frac{\text{Cov}(\widehat{\mathbf{y}},\, \mathbf{y})} {\sqrt{\text{Var}(\widehat{\mathbf{y}})} \, \sqrt{\text{Var}(\mathbf{y})}}.
\end{equation*}

We chose the correlation coefficient since it is widely recognized, easy to interpret, and provides a normalized measure of performance (values lie between $-1$ and $1$). It directly relates to the calibration plots in \autorefs{fig:lmm_alpha_deepdive} and \ref{fig:bhm_alpha_distr_shifts}; a well-calibrated model should exhibit a linear relationship between predictions and observations. For the BTMs, $\widehat{\mathbf{y}}$ represent the conditional mean predictions (more comparable to the SBMs, as opposed to the mean of the posterior predictive distribution, which also includes the modeled noise). As a complementary metric, we used the mean absolute error (MAE), defined as
\vspace{0.25\baselineskip}
\begin{equation*}
    \text{MAE} = \frac{1}{n} \sum_{i=1}^n \left| y_i - \widehat{y}_i \right|.
\end{equation*}

Since the response variables were on a 0--1 scale, this metric can be interpreted as the average percentage point deviation between predicted and observed values.
